% Generated draft: requires mathematical and source review.
\documentclass[11pt]{article}
\usepackage[margin=25mm]{geometry}
\usepackage{fontspec}
\setmainfont{Latin Modern Roman}
\usepackage{amsmath,amssymb,mathtools,unicode-math}
\setmathfont{Latin Modern Math}
\usepackage{enumitem,needspace,xurl,hyperref}
\setlength{\parindent}{0pt}
\setlength{\parskip}{5pt}
\setlength{\emergencystretch}{3em}
\newcommand{\N}{\mathbb{N}}
\newcommand{\Z}{\mathbb{Z}}
\newcommand{\Q}{\mathbb{Q}}
\newcommand{\R}{\mathbb{R}}
\newcommand{\C}{\mathbb{C}}
\newcommand{\F}{\mathbb{F}}
\newcommand{\A}{\mathbb{A}}
\newcommand{\PP}{\mathbb{P}}
\newcommand{\E}{\mathbb{E}}
\newcommand{\eps}{\varepsilon}
\newcommand{\dd}{\,\mathrm{d}}
\begin{document}
\section*{UnsolvedMath ID: 20001053}
\textbf{Source problem code:} catalogue entry
\subsection{Definitions and mathematical statement}
The assigned record is the continuation of a split extraction. The supplied verification gives the preceding heading as ``Conjecture 6.10. (Devos) Under the hypotheses of''; the initial ``Conjecture 6.1'' in this record is the referenced earlier conjecture, not its own number. The appended heading ``6.2 Second \& Kth Neighborhood Conjectures'' belongs to the next section.

An explicit finite labeled-graph interpretation is as follows. Let \(V\) be nonempty, \(n=|V|\), and let \(G_i=(V,E_i)\), \(i=1,\ldots,k\), be simple loopless directed layers, allowing their arc sets to overlap. Labels use \(1<\cdots<k\). A label-increasing path uses strictly increasing chosen labels, and its trivial length-zero path is allowed. A label-increasing rainbow cycle means a directed cycle admitting a starting point and strictly increasing labels along its traversal; no extra inequality from its final label back to its first is imposed. Assume no such cycle. Set
\[
 \delta_i=\min_{v\in V}d_{G_i}^+(v),\qquad
 R(v)=\{w\in V:\text{there is a possibly trivial increasing path from }v\text{ to }w\}.
\]
The aid's most defensible reconstruction of the source claim is
\[
 \frac1n\sum_{v\in V}|R(v)|\ge1+\sum_{i=1}^k\delta_i.
\]
The average is over starting vertices, not a pointwise lower bound for each fixed \(v\), and the minima are global layer minima. Any additional hypotheses imported from Conjecture 6.1 must be checked in the missing surrounding source context.

The source note says the statement can be proved for \(G_i=\operatorname{Cayley}(\Gamma,A_i)\) using Lemma 5.9. Here the natural finite directed Cayley interpretation uses a finite group and arcs \(g\to ga\), \(a\in A_i\); the lemma and its hypotheses are not included, so the note is retained as an attributed special case.

The supplied aid reports the following sharp two-label result for the explicit simple-loopless setting above:
\[
 k=2\quad\Longrightarrow\quad\delta_1+\delta_2\le n-1;
\]
if equality holds, then \(R(v)=V\) for every vertex, so the reconstructed average inequality holds with equality. It does not report a proof of the full average conjecture for every nonextremal instance or every number of labels.

The aid also reports a five-vertex example with \((\delta_1,\delta_2)=(2,1)\) and reachable-set sizes
\[
 (5,3,4,5,5).
\]
Thus one vertex has only three reachable vertices, refuting the plausible pointwise bound \(|R(v)|\ge1+\delta_1+\delta_2=4\), while the average is \(22/5\ge4\). The example's two arc sets are not supplied and cannot be reconstructed uniquely from these counts. An exact layered-DAG double-count identity is also mentioned in the aid, but its formula and auxiliary definitions are absent. These partial results clarify the meaning of the conjecture without resolving it generally.
\subsection{Short English statement}
The source continuation concerns average reachability along increasing labels, with trivial paths counted. The supplied aid proves a reported two-label extremal case and gives reachable-set counts showing that a pointwise interpretation would be false; the full average conjecture remains a separate question.
\subsection{Sources}
\begin{itemize}
\item Local UnsolvedMath registry, permanent ID \(20001053\), source code \texttt{catalogue entry}; \texttt{aim-combinatorics-notes.json}, \texttt{notes[177]}, tag \texttt{conjecture}.
\item AIM workshop notes, ``The Caccetta--H\"aggkvist conjecture,'' split Conjecture 6.10, attributed in the preceding fragment to ``Devos,'' referring to Conjecture 6.1: \url{https://aimath.org/WWN/caccetta/caccetta.pdf}. Statement status \texttt{exact}; the supplied verification refers to original PDF/TeX context. No PDF was accessed for this draft.
\item Source Lemma 5.9, cited for the Cayley-layer case, without the lemma statement or external bibliography. No original remarks or literature field are supplied.
\item Research attempt 1: \texttt{partial\_result}, problem status at run \texttt{open}, rigor label \texttt{proved}, novelty confidence low; registry status \texttt{partially\_solved}, classification \texttt{PARTIAL-PROGRESS}, suggested difficulty L4. Created and updated 14 August 2026. This machine-generated aid is not peer reviewed; independently verify before citation.
\item Ulam AI / UnsolvedMath Contributors; CC BY 4.0.
\end{itemize}
\end{document}
